\section{Introduction}

Transformers owe much of their success to exact token-to-token attention, and much of their
cost to it as well: the attention matrix and the key--value cache grow linearly in sequence
length, so long-context inference pays memory proportional to the number of tokens already
seen \citep{vaswani2017attention,dao2022flashattention}. Recent frontier systems attack this
cost with compressed or sparse attention and with aggressive KV-cache compression
\citep{kimi2026,deepseek2026}. A complementary question is architectural: can a model replace
the pairwise interaction matrix with a \emph{fixed-size} recurrent state that nevertheless
supports several distinct kinds of memory?

SigmaStorm is a concrete answer in code and measurement. One block advances a token-local
latent once:
\begin{equation}
  \label{eq:block}
  \begin{aligned}
    u &= \rms(x),\\
    d_i &= \mathrm{op}_i(u) \quad \text{for each enabled } i,\\
    m &= \frac{\sum_i \alpha_i(u)\,d_i}{\sum_i \alpha_i(u)},\\
    \alpha_i(u) &= \sigma(w_i^\top u),\\
    v &= \Proj\!\Big(u + \Delta t(u)\,\big[m - \nabla E(u) + u\big]\Big),\\
    x' &= x + W_o v .
  \end{aligned}
\end{equation}
The gates $\alpha_i$ are independent sigmoids rather than a softmax; each operator output is
individually RMS-normalised before gating; the mixture is normalised by its gate sum so that
opening more gates does not amplify the proposal; and the projection enforces a hard norm
bound. Depth plays the role of transition time, while recurrent state---not the token prefix---carries
information across positions.

The operators are chosen to cover different failure modes of sequence modeling:
a short convolution for local order, a gated delta rule for key$\to$value associations, a
slot memory for long-horizon prototypes, a ring of pooled chunk summaries for a global path
that neither decays exponentially nor costs $T^2$, optional window attention for precise
local copying, and an energy term that stabilises the latent trajectory. \Cref{tab:operators}
lists their cost and state. Every state tensor has a shape independent of $T$.

\paragraph{Scope and evidence policy.}
This report is deliberately a prototype record. It separates three things that are often
conflated in architecture papers: (i) mechanisms that are \emph{implemented}, (ii) properties
that are \emph{measured} by a falsification gate, and (iii) capabilities that are
\emph{demonstrated} by a controlled experiment. Appendix~\ref{app:ledger} records earlier
claims that failed to reproduce. The central negative results are reported with the same
prominence as the positive ones.

\subsection{Contributions}

\begin{enumerate}[leftmargin=1.2em]
\item \textbf{An executable transition law with a hard state contract.}
      \Cref{sec:method} gives the operator mixture, energy contraction, bounded projection,
      and the exact chunk-carrying rules. The same code runs a 2M-parameter mechanism model
      and a 21M-parameter language model.
\item \textbf{Exact chunk invariance for every operator.}
      The delta rule, slot memory (including its write accumulator), and summary ring
      (including its partial-chunk accumulator) are carried in the state. A gate splits a
      sequence at position 50---inside a 64-token slot block and a 32-token summary block---and
      matches the full forward to $1.6\times10^{-6}$ (\Cref{tab:gates}).
\item \textbf{A multiscale decay spectrum.}
      Each delta-rule head starts at its own log-spaced rate. With the spectrum frozen, recall
      at distances $\geq256$ is $1.000$, versus $0.297$ for a single frozen rate
      (\Cref{sec:results-gap-a}).
\item \textbf{A gated, content-addressed slot read.}
      A novelty dead-zone on writes, a separate read temperature, and a learned
      match-confidence gate turn a diffuse softmax average into an address. Presence recall
      moves from chance to $0.977/1.000/0.953$ over three seeds, and zeroing the memory now
      costs $+2.26$ nats instead of $+0.15$ (\Cref{sec:results-gap-d}).
\item \textbf{No-regression controls.}
      The gated read leaves MQAR ($0.990/0.987$ vs.\ $0.988/0.979$), aggregate recall, and a
      real-data next-token loss ($3.0826$ vs.\ $3.0760$ nats) within seed noise.
\item \textbf{An honest capability map.}
      A 21M model writes coherent TinyStories prose (held-out PPL $3.498$), while state
      expansion, chunk-summary benefit, and 256K-context training remain unproven at this
      scale (\Cref{tab:claims}).
\end{enumerate}

\subsection{Claims and non-claims}

\textbf{We claim}: the mechanisms above are correct under their equivalence tests; the decay
spectrum and the gated slot read have controlled evidence on synthetic recall; a small
language model trains and generates; state size is independent of context.

\textbf{We do not claim}: that SigmaStorm matches or beats transformers at language modeling,
that 256K-token context has been trained or evaluated, that state expansion increases
capacity, that chunk summaries reliably improve quality, or that any result here extrapolates
to frontier scale. The largest model in this report is 21M parameters, trained on 348M tokens
of a synthetic children's corpus; the timing measurements are single-machine CPU forward
passes, not production throughput.
